% Pattern: class/tac. The class of Theory and Applications of Categories
% defines \newtheorem again: each theorem-like environment is a subsection,
% numbered with the subsections, whatever counter the declaration names.
% Seen in the gold pilot (set X).
% want lemma-1 statement/lemma 1.1 status=none
% want proposition-1 statement/proposition 1.2 status=none
% want definition-1 definition 1.4
% want lemma-2 statement/lemma 2.1 status=none
\documentclass{tac}
\newtheorem{lemma}{Lemma}[section]
\newtheorem{proposition}{Proposition}[section]
\newtheorem{definition}{Definition}[section]
\begin{document}
\section{Braidings}
\begin{lemma} Every braiding is invertible. \end{lemma}
\begin{proposition} Every symmetry is a braiding. \end{proposition}
\subsection{Examples}
The vector spaces give an example.
\begin{definition} A braided category is a category with a braiding. \end{definition}
\section{Centres}
\begin{lemma} Every centre is braided. \end{lemma}
\end{document}
